\section{Uvod}\label{sec:uvod}

Izum mehurčne komore pripada Donald A. Glaserju, ki je leta 1952 odkril, da bo pregreta tekočina dietil etra pri temperaturi \( 130^{\degree} \mathrm{C} \) in tlaku 20 atmosfer vzkipnila ob prisotnosti \( ^{60}Co \)~\cite{glaser_effects_1952}. Za svoje izum je leta 1960 prejel Nobelovo nagrado. 

Mehurčne komore vsebujejo pregreto tekočino in, če vpadni delec pri sipanju odda dovolj energije elektronu, se pojavijo mehurčku in tekočina začne vreti~\cite{bressler_buffer-free_2019}. Fotoaparati z dovolj hitrim zajemanjem slik so lahko fotografirale sled mehurčkov, ki jih je za seboj puščal vpadni delec.

Po vretju tekočine se komori poviša tlak, plin se kondenzira, tlak se ponovno zniža in postopek se lahko ponovi~\cite{bressler_buffer-free_2019}. 

Na začetku so se mehurčne komore kot so Gargamelle in Velika evropska mehurčna komora\ (ang. \textit{Big European Bubble Chamber} ali BEBC) uporabljale za sledenje delcem v trkih z visokoenergijskimi delci\cite{noauthor_archives_nodate}. Komore so bile obdane s superprevodnimi magneti, ki so lahko proizvedli magnetno polje do \( 3.5 \mathrm{T} \)~\cite{\BEBC}. Nabiti delci so se, odvisno od predznaka naboja
